\section{When the master comes home}

Jesus' disciples know what a porter does. Someone has to hear the knock, recognize the person outside and open the door. What they want to know is when to expect the knock. They have been sitting opposite the Temple, asking Jesus when the things he has foretold will happen. He ends his answer with a man who leaves home. The master gives his servants authority, entrusts work to them and tells the porter to watch. He gives them nothing from which to calculate the length of his journey. Jesus then turns from the servants to those listening to him. Everyone who hears him is included in the command with which he ends the story (Mark 13:3--4, 33--37).

At the beginning of Advent those words come from someone the Church already knows. At the altar the priest remembers the Son taking our flesh and looks forward to seeing him in glory. Much can happen between the two comings. People who have been enriched by Christ can quarrel over who taught them the faith. A man who once took what belonged to someone else can learn to give to his neighbor. Each of these changes poses a question about the household that expects to see Christ again.

Isaiah's people seem scarcely able to lift a hand. They confess that nobody has stirred himself to lay hold of God. Yet they are speaking to him, and by the end of the prayer they have called themselves the work of his hands. How did a people so estranged come to ask him for help? Jerome listens to the accusation mingled with their plea. Augustine notices, in the psalm that follows Isaiah, precisely who is being asked to turn the people back. Together those observations help explain why Catholics who will shortly offer bread and wine begin by asking even for the will to do what is right.

Paul writes to people in quite a different position. Some have so much to say when the church gathers that one can scarcely be heard over another. Paul thanks God for the gifts of speech and knowledge among them. How can a church so richly provided for still need help to stand firm? Chrysostom attends to the things Paul tells the Corinthians immediately \emph{after} the thanksgiving. Bede, explaining the Gospel, distinguishes the porter's care for the whole house from the charge that falls to each person in it. A gift in one servant's hands can be of great use to another. Paul has reason to worry when those gifts become occasions for members of one church to separate into parties.

The ordinary Sunday response considered here is Psalm 80, the prayer for the vine. The entrance and Communion antiphons are those printed in the Missal. Where a congregation sings one of the common Advent psalms, other verses follow Isaiah. Those possibilities, and the optional blessing of a wreath during Mass, are described in the appendix on liturgical resolution. The Scriptural Date and Location appendix (\ref{sec:date-location}) identifies the passages and their settings.
